% Sections 1-3 and 5-9 (Introduction ... Algorithms). Shared by both templates.

\section{Introduction}
A financial exchange is, at its core, a large concurrent system. Thousands of independent
participants form intentions at the same time, but their orders must be serialised into a single,
strictly ordered book per traded instrument before any trade can occur. Agent-based models (ABMs)
of limit order book (LOB) markets reproduce this structure in software and are widely used to study
price formation, liquidity, and the statistical regularities of returns~\cite{cont2001,preis2006}.
Their main practical limitation is computational: fidelity improves with the number of simulated
agents and instruments, yet mainstream market simulators execute as single-threaded discrete-event
loops~\cite{byrd2020abides} and therefore scale poorly with either.

This workload is interesting for a parallel computer architecture study because it is not
uniformly parallel. The agent decision phase is massively data-parallel and maps naturally onto a
GPU, since every agent independently reads a shared view of the book and emits at most one order.
The matching phase is the opposite: price--time priority imposes a total order on events within an
instrument, so matching for a single instrument resists parallelisation almost entirely and can
only be distributed across instruments. The simulation therefore contains an irreducible serial
core between two parallel regions, and the achievable speedup is governed by the interaction
between them rather than by the efficiency of any single kernel.

This project implements such a simulator with a three-tier heterogeneous mapping: agent decisions
in CUDA, per-instrument matching parallelised across instruments with OpenMP, and instrument sets
distributed across processes with MPI. We treat the resulting performance envelope as the object
of study. The main contributions are:
\begin{itemize}
\item A deterministic LOB simulator whose serial, OpenMP, CUDA and MPI back ends produce
  bit-identical trade logs for a given seed, verified by an order-sensitive hash of every trade.
\item Measurements of the CPU--GPU boundary cost, the effect of warp divergence and copy-back
  compaction, and the agent count at which heterogeneous execution overtakes a multicore baseline.
\item A strong, weak and per-instrument scaling study with Amdahl, Gustafson, Karp--Flatt and
  roofline analysis, showing that order matching is the bottleneck on every back end.
\item An open-source implementation with one-command setup and demo scripts, so every result in
  this report can be reproduced.
\end{itemize}

The rest of the report is organised as follows. Section~\ref{sec:lit} reviews related work,
Sections~\ref{sec:problem} and~\ref{sec:obj} state the problem and objectives,
Sections~\ref{sec:method}--\ref{sec:algo} describe the method, architecture and algorithms,
Sections~\ref{sec:setup}--\ref{sec:perf} present the experiments, results and analysis,
Section~\ref{sec:compare} compares the work with existing methods, and
Sections~\ref{sec:concl}--\ref{sec:future} conclude.

% Section 2: Literature Review. Shared by both templates.
\section{Literature Review}\label{sec:lit}

\subsection{Foundations}
Gode and Sunder showed that agents with no intelligence beyond a budget constraint already produce
near-efficient allocations in a double auction, so much of market behaviour comes from the trading
mechanism itself~\cite{gode1993}. Chiarella and Iori added fundamentalist, chartist and noise
traders interacting through an explicit order book~\cite{chiarella2002}; Preis et al.\ showed that a
simple multi-agent order book with price--time priority reproduces fat tails once order placement
depth depends on the trend~\cite{preis2006}; and Farmer et al.\ found that zero-intelligence order
flow explains much of the observed spread and price impact~\cite{farmer2005}. Cont's stylised facts
(heavy tails, volatility clustering and no linear autocorrelation) are the standard checklist for
validating such models~\cite{cont2001}. On the computing side, Lysenko and D'Souza built the first
largely GPU-resident ABM framework~\cite{lysenko2008}, Fujimoto surveyed parallel discrete-event
simulation~\cite{fujimoto1990}, and Amdahl~\cite{amdahl1967}, Gustafson~\cite{gustafson1988},
Karp and Flatt~\cite{karp1990} and Williams et al.~\cite{williams2009} give the performance models
used in this report.

\subsection{Recent agent-based market simulators}
\textbf{ABIDES-Gym}~\cite{amrouni2021} (ACM ICAIF 2021) wraps ABIDES~\cite{byrd2020abides}, a
message-driven discrete-event simulator with a NASDAQ-like exchange and price--time priority, as an
OpenAI Gym environment for reinforcement learning. It offers high fidelity and realistic latency,
but its kernel is a single-threaded event loop and the paper reports no run times; a later study
cites ABIDES at about 29{,}000 messages per second~\cite{mascioli2024}.

\textbf{JAX-LOB}~\cite{frey2023} (ACM ICAIF 2023) is the first GPU-accelerated LOB simulator. Each
book is a fixed-capacity array, and matching runs inside a JAX loop. The authors note that
processing messages within one book is inherently serial, so all parallelism comes from running
thousands of independent books with \texttt{vmap}: with 1000 books, a message costs about
2.6\,$\mu$s per book on an RTX 2080 Ti, and RL training is about 7$\times$ faster. Its limits are
pre-allocated capacity, a cost set by the slowest branch, and replayed order flow that does not
react to the agent. It does not measure the end-to-end serial fraction or combine GPU agents with
CPU matching.

\textbf{PyMarketSim}~\cite{mascioli2024} (ACM ICAIF 2024) is a Python discrete-time simulator whose
tick (observe, act, clear, advance) matches ours. Its order book takes 1.9--4.7\,$\mu$s per
operation from $10^3$ to $10^6$ orders, about 300{,}000 book operations or 15{,}000 agent actions
per second, but it is serial, simulates one asset and uses 25 agents in its experiments.

\subsection{GPU and high-performance ABM frameworks}
\textbf{FLAME GPU 2}~\cite{richmond2023} (Software: Practice and Experience 2023) compiles agent
functions to CUDA at run time, stores agents as structures of arrays, and passes messages through
spatial, array or bucket structures. It runs up to 200 million agents on a V100 and 1M agents in
about 3\,ms per step. Its performance depends on spatial or keyed locality; agents competing for
one resource need repeated bidding sub-models, atomics are non-deterministic, and one model is
limited to one GPU.

\textbf{Agents.jl}~\cite{datseris2024} (SIMULATION 2024, online 2022) is a Julia ABM library that
runs single models about 2--125$\times$ faster than MASON, NetLogo and Mesa in the authors'
benchmarks. It does not parallelise within a model, because agents mutate shared state, and has no
GPU support: the same shared-state conflict that an order book creates.

\subsection{Matching engines and hardware-accelerated trading}
\textbf{CoinTossX}~\cite{jericevich2022} (SoftwareX 2022) is an open-source Java matching engine
modelled on the Johannesburg Stock Exchange, with price--time priority and messaging over UDP. It
reaches 106--248\,ns 90th-percentile matching latency on a 16-core server, but throughput falls from
2{,}287 to 621--989 orders per second as client--stock pairs are added, and it is engineered for
latency rather than for parallel scaling.

\textbf{LightTrader}~\cite{yoo2023} (IEEE HPCA 2023) is a heterogeneous high-frequency trading system:
an FPGA parses market data and keeps the top of a local order book, while custom deep-learning
accelerators run price-prediction models, with a scheduler that picks batch size and voltage and
frequency under a power budget. It is 13.9$\times$ faster than a V100-based system for AI
processing, but scaling saturates as accelerators share the power budget. It accelerates the
\emph{decision} stage for one trader, the analogue of our GPU phase.

\subsection{Distributed ABM and reproducible parallel random numbers}
\textbf{Repast4Py}~\cite{collier2022} (Winter Simulation Conference 2022) distributes Python ABMs
with MPI, partitioning space across ranks with ghost regions. A 3.6M-agent random walk speeds up
4.39$\times$ from 36 to 288 ranks (efficiency 0.94 to 0.55), while a network model reaches only
1.75$\times$ (efficiency 0.22): communication dominates at higher rank counts. Whether results are
identical across rank counts is not addressed.

\textbf{SIMCoV-GPU}~\cite{leyba2024} (ACM HPDC 2024) ports a spatial infection ABM to multiple GPUs
with CUDA and UPC++, reaching 11.9$\times$ over 512 CPU cores with 16 GPUs. The GPU wins only on
large enough problems (at 64 GPUs on a small grid it is 0.85$\times$), and because CPU and GPU
random number generators differ, results agree only statistically (within about 1\%), not bit for
bit.

\textbf{OpenRAND}~\cite{khan2024} (SoftwareX 2024) provides counter-based generators (Philox,
Threefry and others) keyed by a logical entity rather than a thread, so streams give the same
output on any thread count or device; it is 1.8$\times$ faster than cuRAND in a particle
benchmark. It supplies reproducible random numbers but not reproducibility of a whole simulation,
which also needs deterministic arithmetic and event ordering.

\subsection{Summary, research gaps and motivation}
Table~\ref{tab:lit} summarises the reviewed work. Four gaps motivate this project:
\begin{enumerate}
\item \textbf{Market simulators are built for fidelity, not parallel characterisation.}
  ABIDES-Gym, PyMarketSim and CoinTossX run essentially serially; JAX-LOB parallelises only across
  independent books. None splits the tick into a parallel decision phase and a serial matching
  phase and relates the measured speedup to Amdahl's law.
\item \textbf{The serial matching core is acknowledged but not measured.} No study reports what
  share of tick time matching takes as agents grow, or the gap between kernel and end-to-end
  speedup.
\item \textbf{GPU and distributed ABM frameworks rely on locality.} FLAME GPU 2, Repast4Py and
  SIMCoV-GPU exploit spatial neighbourhoods and halos; Agents.jl avoids in-model parallelism. An
  order book, where every agent writes into one strictly ordered structure per instrument, fits
  none of them.
\item \textbf{Bit-identical heterogeneous results are rare.} SIMCoV-GPU accepts statistical
  agreement and FLAME GPU 2 has non-deterministic atomics; OpenRAND makes only the random numbers
  reproducible. No reviewed work shows a bit-identical event log across CUDA, OpenMP and MPI.
\end{enumerate}
We therefore build a simulator that partitions by instrument, keeps the GPU for decisions only,
produces identical trades on every back end, and measures where its speedup stops and why.

\begin{widetable}
\centering
\caption{Comparison of recent related work (2021--2024).}
\label{tab:lit}
\footnotesize
\begin{tabularx}{\textwidth}{@{}p{2.3cm} X X X p{2.0cm}@{}}
\toprule
Work & Approach / platform & Advantages & Limitations & Key result\\
\midrule
ABIDES-Gym~\cite{amrouni2021} & discrete-event kernel, Python, serial & high fidelity, latency model, RL interface & single-threaded; no run times reported & --\\
JAX-LOB~\cite{frey2023} & JAX on GPU, \texttt{vmap} over books & thousands of books in parallel; RL on GPU & serial within a book; fixed capacity & 2.6\,$\mu$s/msg per book; 7$\times$ RL\\
PyMarketSim~\cite{mascioli2024} & Python discrete-time tick & simple, RL and game-theory tools & serial; one asset; 25 agents & 300k book ops/s\\
FLAME GPU 2~\cite{richmond2023} & CUDA, SoA, spatial messages & 200M agents on one GPU & needs locality; non-det.\ atomics; 1 GPU per model & 1M agents in 3\,ms/step\\
Agents.jl~\cite{datseris2024} & Julia, CPU & concise and fast single models & no in-model parallelism; no GPU & 2--125$\times$ vs other ABM tools\\
CoinTossX~\cite{jericevich2022} & Java engine, UDP & sub-$\mu$s matching latency & throughput falls with instruments & 106--248\,ns latency\\
LightTrader~\cite{yoo2023} & FPGA + DL accelerators & fast decision stage, energy-aware & power-limited scaling; one trader & 13.9$\times$ vs V100 system\\
Repast4Py~\cite{collier2022} & Python + MPI, spatial partitions & HPC scale from Python & communication-bound at high rank counts & 4.39$\times$, 36$\to$288 ranks\\
SIMCoV-GPU~\cite{leyba2024} & CUDA + UPC++, multi-GPU & 11.9$\times$ over 512 cores & GPU wins only at large sizes; not bitwise reproducible & 11.9$\times$ (16 GPUs)\\
OpenRAND~\cite{khan2024} & counter-based RNG library & reproducible on any back end & RNG only, not whole simulation & 1.8$\times$ vs cuRAND\\
\bottomrule
\end{tabularx}
\end{widetable}


\section{Problem Statement}\label{sec:problem}
Given a population of $N$ heterogeneous trading agents and $M$ instruments, each with its own
limit order book, simulate $T$ discrete ticks in which (i) every agent observes a snapshot of its
instrument's book and may submit, cancel or modify one order, (ii) each instrument's orders are
matched under price--time priority, and (iii) a market-wide index is published for the next tick.
The problem is to map this pipeline onto heterogeneous hardware (a GPU, a multicore CPU and
multiple processes) such that:
\begin{enumerate}
\item the parallel versions produce \emph{exactly} the same trades as a serial reference for the
  same random seed, so that correctness can be checked and experiments repeated;
\item the speedup over the serial reference is measured and explained, including the serial
  fraction imposed by price--time priority, the cost of moving data across the CPU--GPU boundary,
  and the cost of synchronising MPI ranks at every tick; and
\item the agent and instrument counts at which the heterogeneous mapping beats an optimised
  multicore CPU implementation are identified.
\end{enumerate}
The simulated market must also remain realistic, reproducing the standard stylised facts of
financial returns~\cite{cont2001}.

\section{Objectives}\label{sec:obj}
The eight objectives from the project synopsis, and where each is addressed, are:
\begin{enumerate}
\item \textbf{Serial matching engine.} Implement a correct serial LOB engine supporting add,
  cancel, modify and market orders under price--time priority, with a unit test suite as the
  correctness reference (Sections~\ref{sec:method}, \ref{sec:algo}, \ref{sec:res-correct}).
\item \textbf{Agents and serial baseline.} Implement noise, momentum, mean-reverting and
  market-making agents and a fully serial simulator as the speedup baseline
  (Section~\ref{sec:method}).
\item \textbf{CUDA agent phase.} Port agent decisions to CUDA with agent-type sorting to reduce warp
  divergence, coalesced access, and counter-based random streams for determinism
  (Sections~\ref{sec:algo}, \ref{sec:res-gpu}).
\item \textbf{OpenMP matching.} Parallelise matching across instruments and evaluate thread scaling,
  false sharing and memory contention (Section~\ref{sec:res-cpu}).
\item \textbf{MPI distribution.} Distribute instruments across ranks with a per-tick index reduction
  and quantify communication overhead (Section~\ref{sec:res-cpu}).
\item \textbf{Validation.} Compare trade logs exactly across back ends and check the stylised facts
  (Sections~\ref{sec:res-correct}, \ref{sec:res-facts}).
\item \textbf{Scaling study.} Strong and weak scaling, an Amdahl fit of the serial fraction,
  roofline placement and Nsight metrics (Sections~\ref{sec:results}, \ref{sec:perf}).
\item \textbf{CPU vs GPU threshold.} Identify where the heterogeneous mapping overtakes the
  multicore CPU and explain why (Sections~\ref{sec:res-gpu}, \ref{sec:perf}).
\end{enumerate}

\section{Proposed Methodology}\label{sec:method}

\subsection{Barrier-synchronised tick}
The simulator advances in discrete ticks. Following the conservative synchronisation family of
parallel discrete-event simulation~\cite{fujimoto1990}, every tick has three phases separated by
barriers:
\begin{enumerate}
\item \textbf{Decide} (parallel over agents): each agent reads its instrument's snapshot (best bid,
  best ask, last price, trend) and the market index, and emits at most one order intent.
\item \textbf{Match} (parallel over instruments, serial inside each): the intents of each instrument
  are applied to its book in a deterministic arrival order, producing trades.
\item \textbf{Publish} (reduction over instruments): snapshots for the next tick are rebuilt and the
  market index (the sum of mid prices) is reduced, across ranks with \texttt{MPI\_Allreduce}.
\end{enumerate}
This trades some concurrency for determinism and auditability; the cost of the barrier is one of
the quantities measured.

\subsection{Order book}
Each instrument has a fixed grid of 4096 integer price levels with one FIFO queue per level. Orders
live in a node pool and are doubly linked into their level, with an id-to-node hash map, so add,
cancel and modify are $O(1)$ except for moving the best-price pointer when a level empties. Limit
orders match against the opposite side at the resting (maker) price and the remainder rests;
market orders match and drop any unfilled remainder. A modify keeps time priority only for a
quantity decrease at the same price. Resting orders expire after a time-to-live (100 ticks by
default), which keeps the book bounded.

\subsection{Agent population}
Four strategies, after the heterogeneous populations of~\cite{chiarella2002,preis2006}, are written
once as \texttt{\_\_host\_\_ \_\_device\_\_} functions so that the CPU loops and the CUDA kernel run
literally the same code (Table~\ref{tab:agents}). Agents are stateless functions of (agent id,
tick, seed, snapshot): they do not track inventory, which keeps the GPU phase free of per-tick
feedback from the matching engine.

\begin{table}[t]
\centering
\caption{Agent population (default mix).}
\label{tab:agents}
\footnotesize
\begin{tabularx}{\columnwidth}{@{}l c c X@{}}
\toprule
Strategy & Share & Acts/tick & Behaviour\\
\midrule
Noise & 65\% & 10\% & random side; mostly limit orders near the reference price\\
Momentum & 15\% & 5\% & market order in the direction of the instrument and index trend\\
Mean-reverting & 10\% & 3\% & limit order back toward the fundamental value\\
Market maker & 10\% & 100\% & re-quotes one side per tick, cancels its quote from two ticks earlier\\
\bottomrule
\end{tabularx}
\end{table}

\subsection{Determinism by design}
Parallel stochastic simulations are usually not repeatable, because random draws and floating-point
reductions depend on thread scheduling. We make every back end produce the identical trade log:
\begin{itemize}
\item \textbf{Counter-based RNG:} Philox4x32-10~\cite{salmon2011} maps (seed, stream, agent, tick)
  to random numbers with no shared state, so thread order cannot change any draw.
\item \textbf{Integer-only decisions:} CPU and GPU floating point can round differently (fused
  multiply--add, math libraries), so decisions use integer arithmetic.
\item \textbf{Deterministic arrival order:} intents are applied in agent-id order, rotated by a
  random offset per (instrument, tick) so that no agent is always first.
\item \textbf{Predictable order ids:} an order's id is $(\textit{tick}+1)\ll 32 \mid \textit{agent}$,
  so a market maker on the GPU can cancel its earlier quote without feedback from the CPU.
\item \textbf{Integer reductions:} the index is a sum of integers, so \texttt{MPI\_Allreduce} gives
  the same result in any reduction order.
\item \textbf{Order-independent compaction:} GPU output slots are reserved with
  \texttt{atomicAdd}; each entry carries its agent index and the host scatters by index.
\end{itemize}
Correctness is checked with an order-sensitive FNV-1a hash over every trade (tick, price,
quantity, buyer, seller, maker and taker ids), combined across instruments and ranks.

\subsection{Parallel mapping}
\begin{itemize}
\item \textbf{CUDA (decide):} one thread per agent, structure-of-arrays layout for coalesced loads,
  agents sorted by strategy so that warps mostly take one branch, snapshots read through the
  read-only cache (\texttt{\_\_ldg}), parameters in constant memory, and compaction of the copy
  back to only the agents that acted (about 20\% per tick).
\item \textbf{OpenMP (decide and match):} \texttt{parallel for schedule(static)} over agents for decisions;
  \texttt{parallel for schedule(dynamic,1)} over instruments for matching, with per-thread counters
  padded to 64-byte cache lines to avoid false sharing.
\item \textbf{MPI (instruments):} each rank owns a contiguous block of instruments and the agents
  trading them; ranks meet once per tick in an integer \texttt{MPI\_Allreduce}.
\end{itemize}

\subsection{Measurement}
The program times every phase of every tick (decide, host-to-device copy, kernel, device-to-host
copy, match, publish, MPI communication) and reports totals; for MPI it reports the slowest rank.
Kernel time uses CUDA events. Speedup is $S(p)=T_1/T_p$ against the serial back end, efficiency
is $E(p)=S(p)/p$, the serial fraction $f$ is fitted to Amdahl's law
$S(p)=1/(f+(1-f)/p)$~\cite{amdahl1967} by least squares, and the Karp--Flatt metric
$e(p)=(1/S-1/p)/(1-1/p)$~\cite{karp1990} is reported for each $p$. Weak scaling is interpreted
with Gustafson's law~\cite{gustafson1988} and the kernel is placed on a roofline~\cite{williams2009}
from Nsight Compute counters.

\section{System Architecture}\label{sec:arch}
Fig.~\ref{fig:arch} shows how the three tiers map onto hardware and Fig.~\ref{fig:flow} shows the
control flow of one simulation run. MPI splits the instruments into blocks, one per rank. Inside a
rank, the decision phase runs on the GPU (or on all CPU cores with OpenMP), and the matching phase
runs on the CPU with one instrument per OpenMP thread at a time. The ranks synchronise once per
tick to reduce the market index. Table~\ref{tab:modules} lists the software modules.

\begin{figure}[t]
\centering
\resizebox{\columnwidth}{!}{%
\begin{tikzpicture}[
  font=\sffamily\small,
  box/.style={draw, rounded corners=2pt, align=center, minimum height=0.8cm, inner sep=3pt},
  gpu/.style={box, fill=green!12},
  cpu/.style={box, fill=blue!10},
  net/.style={box, fill=orange!15},
  arr/.style={-{Latex[length=2mm]}, thick}]
  % Rank 0
  \node[draw, dashed, rounded corners, minimum width=6.0cm, minimum height=5.0cm,
        label={[font=\sffamily\bfseries\small]above:MPI rank 0 (instruments $0..M/2-1$)}] (r0) at (0,0) {};
  \node[gpu, minimum width=5.4cm] (g0) at (0,1.5) {GPU: \textbf{decide} kernel\\one thread per agent, SoA, sorted};
  \node[box, minimum width=5.4cm, fill=gray!10] (c0) at (0,0.2) {compaction (atomicAdd) $\rightarrow$ copy back\\host scatter by agent index};
  \node[cpu, minimum width=5.4cm] (m0) at (0,-1.3) {CPU: \textbf{match} with OpenMP\\thread $k$ $\leftarrow$ instrument $i$ (dynamic)};
  \draw[arr] (g0) -- (c0); \draw[arr] (c0) -- (m0);
  % Rank 1
  \node[draw, dashed, rounded corners, minimum width=6.0cm, minimum height=5.0cm,
        label={[font=\sffamily\bfseries\small]above:MPI rank 1 (instruments $M/2..M-1$)}] (r1) at (6.8,0) {};
  \node[gpu, minimum width=5.4cm] (g1) at (6.8,1.5) {GPU: \textbf{decide} kernel\\one thread per agent, SoA, sorted};
  \node[box, minimum width=5.4cm, fill=gray!10] (c1) at (6.8,0.2) {compaction (atomicAdd) $\rightarrow$ copy back\\host scatter by agent index};
  \node[cpu, minimum width=5.4cm] (m1) at (6.8,-1.3) {CPU: \textbf{match} with OpenMP\\thread $k$ $\leftarrow$ instrument $i$ (dynamic)};
  \draw[arr] (g1) -- (c1); \draw[arr] (c1) -- (m1);
  % Allreduce
  \node[net, minimum width=12.6cm] (ar) at (3.4,-3.4) {\textbf{publish}: rebuild snapshots, \texttt{MPI\_Allreduce} of the integer index (barrier)};
  \draw[arr] (m0) -- (m0|-ar.north); \draw[arr] (m1) -- (m1|-ar.north);
  \draw[arr, dashed] (ar.west) -- ++(-0.6,0) |- node[pos=0.25, left, align=right]{next\\tick} (g0.west);
  \draw[arr, dashed] (ar.east) -- ++(0.6,0) |- (g1.east);
\end{tikzpicture}}
\caption{System architecture: the three-tier heterogeneous mapping (two MPI ranks shown).}
\label{fig:arch}
\end{figure}

\begin{figure}[t]
\centering
\resizebox{0.86\columnwidth}{!}{%
\begin{tikzpicture}[
  font=\sffamily\small, node distance=0.42cm,
  st/.style={draw, rounded corners=8pt, fill=gray!15, minimum width=4.2cm, minimum height=0.7cm, align=center},
  pr/.style={draw, fill=blue!8, minimum width=4.2cm, minimum height=0.7cm, align=center},
  dc/.style={draw, diamond, aspect=2.4, fill=orange!15, align=center, inner sep=1pt},
  arr/.style={-{Latex[length=2mm]}, thick}]
  \node[st] (start) {Start: parse options (N, M, T, seed, back ends)};
  \node[pr, below=of start] (init) {Create agents (sorted by type),\\books, Philox streams; split instruments over ranks};
  \node[pr, below=of init] (snap) {Build snapshot of every book};
  \node[pr, below=of snap, fill=green!12] (dec) {\textbf{Decide}: serial / OpenMP / CUDA\\each agent emits $\leq 1$ intent};
  \node[pr, below=of dec] (arrv) {Order intents per instrument\\(rotated agent-id order)};
  \node[pr, below=of arrv, fill=blue!12] (mat) {\textbf{Match}: one instrument per thread,\\price--time priority, hash trades, expire old orders};
  \node[pr, below=of mat, fill=orange!12] (pub) {\textbf{Publish}: market index by\\\texttt{MPI\_Allreduce} (barrier)};
  \node[dc, below=of pub] (more) {tick $< T$?};
  \node[st, below=of more] (out) {Combine hashes over ranks; print\\trades, fingerprint and phase times};
  \draw[arr] (start) -- (init); \draw[arr] (init) -- (snap); \draw[arr] (snap) -- (dec);
  \draw[arr] (dec) -- (arrv); \draw[arr] (arrv) -- (mat); \draw[arr] (mat) -- (pub);
  \draw[arr] (pub) -- (more); \draw[arr] (more) -- node[right]{no} (out);
  \coordinate (R) at ($(mat.east)+(0.5,0)$);
  \draw[arr] (more.east) -- node[above]{yes} (more.east -| R) |- (snap.east);
\end{tikzpicture}}
\caption{Flowchart of one simulation run.}
\label{fig:flow}
\end{figure}

\begin{table}[t]
\centering
\caption{Software modules (C++17, CUDA, OpenMP, MPI).}
\label{tab:modules}
\footnotesize
\begin{tabularx}{\columnwidth}{@{}l X@{}}
\toprule
File & Role\\
\midrule
\texttt{order\_book.hpp/.cpp} & serial matching engine, price--time priority\\
\texttt{agents.hpp} & four strategies, shared by CPU and GPU\\
\texttt{philox.hpp} & Philox4x32-10 counter-based RNG\\
\texttt{sim.hpp/.cpp} & tick loop, OpenMP matching, MPI reduction, trade hash\\
\texttt{decide\_cuda.cu} & CUDA decision kernel and compaction\\
\texttt{main.cpp} & \texttt{lobsim} command-line program\\
\texttt{test\_main.cpp} & unit tests (259{,}597 checks)\\
\texttt{scripts/}, \texttt{analysis/} & setup, demo, benchmarks, profiling, plots\\
\bottomrule
\end{tabularx}
\end{table}

\section{Algorithms}\label{sec:algo}
Algorithm~\ref{alg:tick} is the top-level tick loop executed by every MPI rank.
Algorithm~\ref{alg:match} matches one incoming limit order under price--time priority; it is the
serial core of the simulation. Algorithm~\ref{alg:kernel} is the CUDA decision kernel.

\begin{algorithm}[t]
\caption{Barrier-synchronised tick loop (per MPI rank)}
\label{alg:tick}
\footnotesize
\begin{algorithmic}[1]
\Require agents $A_r$ and instruments $I_r$ owned by rank $r$; ticks $T$; seed $s$
\For{$t \gets 0$ \textbf{to} $T-1$}
  \State $snap[i] \gets \textsc{Snapshot}(book[i])$ \textbf{for all} $i \in I_r$
  \State \textit{// Phase 1: decide (CUDA kernel or OpenMP parallel for)}
  \ForAll{$a \in A_r$ \textbf{in parallel}}
    \State $intent[a] \gets \textsc{Decide}(a, t, s, snap[inst(a)], index)$
  \EndFor
  \State \textit{// Phase 2: match, one instrument per thread}
  \ForAll{$i \in I_r$ \textbf{in parallel} (\texttt{schedule(dynamic,1)})}
    \State $o \gets \textsc{Rotate}(\text{agents of } i, \textsc{Philox}(s,i,t))$
    \ForAll{$a$ \textbf{in order} $o$ with $intent[a] \neq \emptyset$}
      \State $trades \gets \textsc{Apply}(book[i], intent[a])$ \Comment{Alg.~\ref{alg:match}}
      \State $h[i] \gets \textsc{FNV1a}(h[i], trades)$
    \EndFor
    \State \textsc{ExpireOrders}$(book[i], t)$
  \EndFor
  \State \textit{// Phase 3: publish}
  \State $local \gets \sum_{i\in I_r} mid(book[i])$ \Comment{integer}
  \State $index \gets \textsc{MPI\_Allreduce}(local, \textsc{Sum})$ \Comment{barrier}
\EndFor
\State \Return $\textsc{MPI\_Allreduce}(h, \textsc{BXor})$ \Comment{per-instrument hashes}
\end{algorithmic}
\end{algorithm}

\begin{algorithm}[t]
\caption{Match a limit buy order (price--time priority)}
\label{alg:match}
\footnotesize
\begin{algorithmic}[1]
\Require book with ask levels $L[p]$ (FIFO queues), best ask $p_a$; order $(id, p, q)$
\While{$q > 0$ \textbf{and} $p_a \leq p$}
  \State $m \gets \textsc{Front}(L[p_a])$ \Comment{oldest order at best price}
  \State $x \gets \min(q, m.q)$
  \State \textsc{EmitTrade}$(p_a, x, \text{maker}=m, \text{taker}=id)$
  \State $q \gets q - x$; \ $m.q \gets m.q - x$
  \If{$m.q = 0$}
    \State remove $m$ from $L[p_a]$ and from the id map
    \If{$L[p_a]$ is empty} \State $p_a \gets$ next non-empty ask level \EndIf
  \EndIf
\EndWhile
\If{$q > 0$} \State append $(id, p, q)$ to the back of bid level $p$; update best bid \EndIf
\end{algorithmic}
\end{algorithm}

\begin{algorithm}[t]
\caption{CUDA decision kernel with compaction}
\label{alg:kernel}
\footnotesize
\begin{algorithmic}[1]
\Require SoA arrays $type[], inst[]$ (sorted by type); snapshots in global memory; params in constant memory
\State $a \gets blockIdx.x \cdot blockDim.x + threadIdx.x$
\If{$a \geq N$} \Return \EndIf
\State $sn \gets \texttt{\_\_ldg}(\&snap[inst[a]])$ \Comment{read-only cache}
\State $(r_0..r_3) \gets \textsc{Philox4x32}(s, a, t)$ \Comment{no shared RNG state}
\State $o \gets \textsc{Decide}_{type[a]}(sn, index, r)$ \Comment{integer only; warps mostly uniform}
\If{$o \neq \emptyset$}
  \State $k \gets \texttt{atomicAdd}(\&count, 1)$
  \State $out[k] \gets (a, o)$ \Comment{slot order varies; agent index makes it irrelevant}
\EndIf
\State \textit{// Host: copy back $out[0..count-1]$ only, then $intent[out[k].a] \gets out[k].o$}
\end{algorithmic}
\end{algorithm}

\textbf{Complexity.} Decisions cost $O(N)$ work per tick with $O(N/p)$ depth on $p$ processors.
Each order costs $O(1)$ to add or cancel plus $O(k)$ for $k$ fills; with $n_i$ orders on instrument
$i$, matching costs $O(\sum_i n_i)$ work but its depth is $O(\max_i n_i)$, because the orders of
one instrument are processed sequentially. With a single instrument the depth equals the work,
which is the source of the serial fraction measured in Section~\ref{sec:perf}.
